\documentclass[english, parskip=full]{scrartcl}
\usepackage[utf8]{inputenc}  % Kodierung der Datei
\usepackage[english]{babel}  % Sprache auf Deutsch 
\usepackage{color}
\usepackage{graphicx, gensymb}
\usepackage{amsfonts, cancel, amssymb}
\usepackage{amsmath, amsthm, array, esint}
\usepackage{booktabs, multirow}  % schönere Tabellen
\usepackage{caption}  % Anpassung der Captions
\usepackage{scrlayer-scrpage}    % Manipulation des Seitenstils
\pagestyle{scrheadings}  % Stil für die Seite setzen
\automark{section}  % Abschnittsnamen als Seitenbeschriftung 
\setheadsepline{.5pt}    % Kopzeile durch Linie abtrennen
\usepackage[hidelinks]{hyperref}  % Links und weitere PDF-Features
\usepackage{typearea, tikz, pgffor, verbatim}
\usetikzlibrary{calc, arrows.meta,arrows, graphs, quotes}
\usepackage[cache=false, outputdir=build]{minted}
\usemintedstyle{vs}
\usepackage{placeins} % provides \FloatBarrier

\definecolor{bg}{rgb}{0.9,0.9,0.9}
\newcommand\numberthis{\addtocounter{equation}{1}\tag{\theequation}}
\usepackage{svg}
\usepackage{siunitx}
\usepackage{physics}
\usepackage{tensor}
\renewcommand{\bra}{}
\newcommand{\kom}{}
\newcommand{\brak}{}
\renewcommand{\braket}{}
\newcommand{\intx}{}
\newcommand{\intr}{}
\newcommand{\kla}{}
\newcommand{\klam}{}
\newcommand{\klammer}{}
\newcommand{\oper}{}
\newcommand{\tecxt}{}
\newcommand{\Bra}{}
\newcommand{\Ket}{}
\renewcommand{\i}{\text{i}}
\newcommand{\e}{\text{e}}
\newcommand*{\Fourier}{\textsc{Fourier}\ }
\newcommand*{\F}{\mathcal{F}}
\newcommand*{\sinc}{\text{sinc\,}}
\newcommand{\conv}{\mathop{\scalebox{3.5}{\raisebox{-0.38ex}{\makebox[0.5\width][c]{$\ast$}}}}\limits}

\newcommand*{\titel}{Gravitational effects on light}
\newcommand*{\autor}{Joachim Schwardt}

\hypersetup{pdfauthor={\autor}, pdftitle={\titel}}

%\titlehead{titlehead}
\subject{General Relativity}
\title{\titel}
\author{\autor}
\date{\today}

\begin{document}

%\begin{titlepage}
\maketitle  % Titel setzen
%\tableofcontents  % Inhaltsverzeichnis setzen
%\end{titlepage}

\section{Light bending within Newtonian mechanics}
Before we discuss how the gravitational effects on light can be described in Newtonian mechanics, we want to reiterate the solution to the two-body problem.
However, we will assume that the heavy object with mass $M$ does not move, and merely acts as a gravitational field for the projectile of mass $m$.
This leads to the total energy
\begin{align}
E &= \frac{m}{2} \dot{\vec{r}}\,^2 - \frac{GMm}{|\vec{r}|}.
\end{align}
Here we have $\vec{r} = r\vec{e}_r + z\vec{e}_z$, where $r = \sqrt{x^2 + y^2}$ with $x = r\cos\phi$ and $y=r\sin\phi$ being the Cartesian coordinates.
Note that the angular momentum is given by
\begin{align}
\vec{L} &= \vec{r} \times \vec{p} = m\kla\left( r\vec{e}_r + z\vec{e}_z \right) \times \kla\left( \dot{r}\vec{e}_r + r\dot{\phi}\vec{e}_\phi + \dot{z}\vec{e}_z \right) = \\
&= mr^2\dot{\phi} \vec{e}_z - mr\dot{z}\vec{e}_\phi + mz\dot{r}\vec{e}_\phi - mz r\dot{\phi} \vec{e}_r.
\end{align}
Since there is no force acting along $\vec{e}_z$, we transform into a frame of reference in which $z \equiv 0$.
Then the angular momentum simplifies to 
\begin{align}
\vec{L} &= mr^2\dot{\phi} \vec{e}_z \equiv L\vec{e}_z.
\end{align}
Inserting into the energy equation gives
\begin{align}
E &= \frac{m}{2} \kla\left(  \dot{r}^2 + r^2\dot{\phi}^2 \right) - \frac{GMm}{r} = \frac{m}{2} \dot{r}^2 + \frac{L^2}{2mr^2} - \frac{GMm}{r}.
\end{align}
We may now consider $r(\phi(t))$ instead of $r(t)$, leading to
\begin{align}
E &= \frac{m}{2} \kla\left( \frac{\tecxt\text{d}r(\phi)}{\tecxt\text{d}\phi} \dot{\phi} \right)^2 + \frac{L^2}{2mr^2} - \frac{GMm}{r} = \\
&= \frac{L^2r'(\phi)^2}{2mr^4}  + \frac{L^2}{2mr^2} - \frac{GMm}{r}.
\end{align}
Introducing the new variable $u := \frac{1}{r}$ with $r' = -\frac{u'}{u^2}$ then results in
\begin{align}
E &= \frac{L^2u'(\phi)^2}{2m} + \frac{L^2 u(\phi)^2}{2m} - GMm u(\phi).
\end{align}
Differentiating and dividing by $\frac{L^2}{m u'(\phi)}$ gives
\begin{align}
u'' + u &= \frac{GMm^2}{L^2} =: \gamma. \label{eqn:newtonian trajectory general solution}
\end{align}
This equation has the general solution
\begin{align}
u(\phi) &= \gamma + A\sin\phi + B\cos\phi,
\end{align}
where we now need to relate the coefficients $A,B$ to the initial condition.
The particular setup we want to investigate is sketched in \autoref{figure:sketch newtonian trajectory}.
Note that the initial velocity is taken to be parallel to $\vec{e}_x$, and that the heavy object has a radius of $R$.
\begin{figure}[!ht]
    \centering
    \begin{tikzpicture}[scale=1,cap=round]
        \pgfmathsetmacro{\R}{1.0}
        \pgfmathsetmacro{\r}{0.05}
        \pgfmathsetmacro{\y}{1.4}
        
        \coordinate (o) at (0,0);
        \coordinate (xmax) at (5,0);
        \coordinate (xmin) at (-7,0);
        \coordinate (ymax) at (0,3);
        \coordinate (ymin) at (0,-1.5);
        
        \coordinate (r0) at (-6,\y);
        \coordinate (r0right) at (-5,\y);
        
        \coordinate (r1) at (0,\y);
        \coordinate (r2) at (\R,0.8*\y);
        \coordinate (r3) at (4.8*\R,-0.4);
        
%        \node[anchor=north east] at (o) {$0$};
        \node[anchor=north] at (xmax) {$x$};
        \node[anchor=west] at (ymax) {$y$};
        
        \draw[fill, color=black!35!white] (o) circle (\R);
        \draw (o) circle (\R);
        \node[below=0.5cm, right] at (0,\R) {$R$};
        \draw[color=black, -{Stealth[length=1.5mm,width=1.5mm]}] (o) to (0,\R);
        \node[below=0.5cm, left] at (o) {$M$};
        
        \draw[fill] (r0) circle (\r);
        \node[below] at (r0) {$m$};
        \node[above] at (r0) {$\vec{r}_0$};
        
        \draw[color=black, thick=0.05cm, -{Stealth[length=1mm,width=1mm]}] (r0) .. controls (r1) and (r2)  .. (r3);
        \draw[color=blue, thick=0.05cm, -{Stealth[length=1.5mm,width=1.5mm]}] (r0) to (r0right); 
        \node[above, color=blue] at (r0right) {$\vec{v}_0$};      
        
        
        \draw[color=black, -{Stealth[length=1.5mm,width=1.5mm]}] (xmin) to (xmax);
        \draw[color=black, -{Stealth[length=1.5mm,width=1.5mm]}] (ymin) to (ymax);
        
        
        
%        \coordinate (A) at (0,0);
%        \coordinate (B) at (1.5,1);
%        \coordinate (C) at (0,1);
%        \coordinate (D) at (1.5,0);
%        \coordinate (M) at (0.75,0.5);
%        
%        \coordinate (CO) at (-0.25, -0.15);
%        \coordinate (CT) at (-0.25, 1.25);
%        \coordinate (CR) at (1.75, -0.15);
%        \coordinate (CM) at (0.75, -0.15);
%        
%        \draw[color=blue, thick] (A) .. controls (M) ..  (D);
%        \draw[color=blue, thick] (B) .. controls (M) ..  (C);
%        \draw[dashed] (A) to (B);
%        \draw[dashed] (C) to (D);
%        
%        \node[anchor=east] at (A) {$\ket{1}$};
%        \node[anchor=west] at (B) {$\ket{1'}$};
%        \node[anchor=east] at (C) {$\ket{2}$};
%        \node[anchor=west] at (D) {$\ket{2'}$};
%        
%        
%        \node[anchor=south,text=blue] at ($(M)+(0,0.2)$) {$\ket{\phi_2}$};
%        \node[anchor=north,text=blue] at ($(M)-(0,0.2)$) {$\ket{\phi_1}$};
%        
%        \draw[color=black, -{Stealth[length=3mm,width=2mm]}] (CO) to (CT);
%        \draw[color=black, -{Stealth[length=3mm,width=2mm]}] (CO) to (CR);
%        
%        \node[anchor=south] at (CT) {$V(z)$};
%        \node[anchor=west] at (CR) {$z$};
%        
%        \draw[color=black] ($(CM)+(0,-0.025)$) to ($(CM)+(0,0.025)$);
%        \node[anchor=north west] at (CM) {$z_c$};
        
    \end{tikzpicture}
    \caption{Sketch of the trajectory of $m$ within the gravitational field of the sphere with mass $M$ and radius $R$.
    Note that we assume the initial velocity to be $\vec{v}_0 = v_0\vec{e}_x$.}
    \label{figure:sketch newtonian trajectory}
\end{figure}\\
Since we are dealing with a central potential, the angular momentum is conserved.
Thus we may compute it at $t=0$, for which we will go back to the Cartesian representation.
Then we find
\begin{align}
\vec{L} &= m\vec{r} \times \dot{\vec{r}} = m\kla\left( x\vec{e}_x + y\vec{e}_y \right) \times \kla\left( \dot{x}\vec{e}_x + \dot{y}\vec{e}_y \right) = m\kla\left( x\dot{y} - y\dot{x} \right)\vec{e}_z.
\end{align}
Since $\dot{y}(t=0) \overset{!}{=} 0$ and $\dot{x}(t=0) = v_0$, this simplifies to 
\begin{align}
L &= -my_0v_0.
\end{align}
From \autoref{eqn:newtonian trajectory general solution} we can also derive
\begin{align}
\frac{1}{x(t)} &= \frac{u(\phi(t))}{\cos\phi(t)} = B + A\tan\phi(t) + \gamma\sec\phi (t), \\
\frac{1}{y(t)} &= \frac{u(\phi(t))}{\sin\phi(t)} = A + B\cot\phi(t) + \gamma\csc\phi(t).
\end{align}
Since $x_0 = -|x_0|$ and $y_0 > 0$, we find $\phi_0 = \pi - \arctan\left\vert \frac{y_0}{x_0} \right\vert$.
Using $\frac{\tecxt\text{d}}{\tecxt\text{d}t} \frac{1}{x} = -\frac{\dot{x}}{x^2}$, we can express the initial velocities via
%Then the above can be rewritten as
%\begin{align}
%\frac{1}{x_0} &= B - A\frac{y_0}{|x_0|} - \gamma \sqrt{1 + \frac{y_0^2}{x_0^2}},\\
%\frac{1}{y_0} &= A - B\frac{|x_0|}{y_0} + \gamma \sqrt{1 + \frac{x_0^2}{y_0^2}}.
%\end{align}
\begin{align}
\dot{x}(t=0) &= -x_0^2 \frac{\tecxt\text{d}}{\tecxt\text{d}t} \frac{1}{x(t)} = -x_0^2 \dot{\phi}(t) \sec^2\phi(t) \kla\left( A + \gamma\sin\phi(t) \right)  \Big\vert_{t=0}, \\
\dot{y}(t=0) &= -y_0^2 \frac{\tecxt\text{d}}{\tecxt\text{d}t} \frac{1}{y(t)} = y_0^2 \dot{\phi}(t) \csc^2\phi(t) \kla\left( B + \gamma\cos\phi(t) \right)  \Big\vert_{t=0}.
\end{align}
Note that $\dot{\phi}(t=0) = \frac{L}{mr_0^2} = \frac{L}{m(x_0^2+y_0^2)}$.
Then we can solve for the desired coefficients
\begin{align}
A &= -\gamma \sin\phi_0 - \frac{v_0}{x_0^2\sec^2\phi_0} \frac{m(x_0^2+y_0^2)}{L} = \\
&= -\gamma \frac{1}{\sqrt{1 + \frac{x_0^2}{y_0^2}}} - \frac{v_0}{x_0^2 \kla\left( 1 + \frac{y_0^2}{x_0^2} \right)} \frac{m(x_0^2 + y_0^2)}{L} = \\
&= -\frac{\gamma y_0}{r_0} - \frac{mv_0}{L},\\
B &= -\gamma\cos\phi_0 = \frac{\gamma}{\sqrt{1 + \frac{y_0^2}{x_0^2}}} = \frac{\gamma |x_0|}{r_0}.
\end{align}
The solution for our particular setup can then be written as
\begin{align}
u(\phi) &= \gamma - \kla\left( \frac{\gamma y_0}{r_0} + \frac{mv_0}{L} \right) \sin\phi + \frac{\gamma |x_0|}{r_0} \cos\phi.
\end{align}
In this form we may now safely take the limit $x_0 \rightarrow -\infty$, which corresponds to the projectile $m$ entering the system as a quasi-free particle.
Then we find
\begin{align}
u(\phi) &= \gamma - \frac{mv_0}{L}\sin\phi + \gamma \cos\phi = \\
&= 2\gamma \cos^2\frac{\phi}{2} - 2\frac{mv_0}{L} \sin\frac{\phi}{2} \cos\frac{\phi}{2}.
\end{align}
To understand by how much the projectile is deflected, we need to figure out the difference between the incoming and outgoing angles $\phi_\tecxt\text{out} - \phi_\tecxt\text{in}$.
However, without any influence we would expect this difference to amount to $180^\circ \equiv \pi$, so we define the deviation as
\begin{align}
\Delta \phi &:= \pi + \phi_\tecxt\text{in} - \phi_\tecxt\text{out} \equiv -\phi_\tecxt\text{out}\ \tecxt\text{mod}\ 2\pi.
\end{align}
Note that by ``outgoing'' we refer to $r(\phi) \rightarrow \infty$ as $\phi\rightarrow\phi_\tecxt\text{out}$.
We can find this angle by solving $u(\phi) = 0$.
This leads to
\begin{align}
\tan\frac{\phi}{2} &= \frac{\gamma L}{mv_0}.
\end{align}
Therefore we finally arrive at the result
\begin{align}
\Delta\phi &= -2\arctan \frac{\gamma L}{mv_0} = -2\arctan
\frac{GMm}{Lv_0} = 2\arctan\frac{GM}{y_0v_0^2}. \label{eqn:newtonian trajectory deflection angle}
\end{align}
Notice that the projectile mass $m$ cancelled out.
For $v_0\rightarrow c$ and $y_0 = R$ we then find
\begin{align}
\Delta\phi &= 2\arctan\frac{GM}{Rc^2} \simeq \frac{2GM}{Rc^2}.
\end{align}
The approximation is usually valid for stars.
For the sun, we have $M=\SI{1.989e30}{kg}$ and $R=\SI{6.963e8}{m}$ giving $\Delta\phi = 0.875''$.


%\begin{thebibliography}{9}
%\bibitem{example} description, \url{https://www.exmapleurl}, day.\,month~2021.
%\end{thebibliography}

\end{document}
